\chapter{Design and Implementation}
\label{ch:design}

% BUDGET: 12 pages. Design and implementation are presented together, as the
% submission brief permits: the system was built iteratively, and most of its
% design decisions are only intelligible alongside the failure that forced them.

This chapter presents the design and the implementation together. They are not
separable in this project: almost every structural decision described below was made in
response to a specific observed failure, and presenting the structure without the failure
would misrepresent both how the system came to look the way it does and why it is right that
it does. The chapter therefore proceeds from the architecture outwards, and each mechanism is
introduced with the problem that produced it.

\section{Architecture}
\label{sec:architecture}

The system is a pipeline of seven agents, shown in Figure~\ref{fig:architecture}. Each agent
is a self-contained stage with a published output contract and a validation function that
enforces it. Two orchestrators sit above them, and neither is an agent: a small loop
coordinating the Writer and Reviewer, and a top-level graph running all seven stages
end to end.

\begin{figure}[htbp]
\centering
\begin{tikzpicture}[
  node distance=4mm,
  ag/.style={draw,rounded corners=1pt,align=center,font=\scriptsize,
             minimum height=7mm,minimum width=20mm,fill=black!3},
  lbl/.style={font=\scriptsize\itshape},
  ar/.style={-{Stealth[length=1.4mm]},thick}
]
\node[ag] (lit) {Literature};
\node[ag,right=of lit] (ida) {Interdisc.\\Literature};
\node[ag,right=of ida] (hyp) {Hypothesis};
\node[ag,right=of hyp] (plan) {Experiment\\Planner};
\node[ag,below=8mm of plan] (code) {Coder};
\node[ag,left=of code] (wri) {Writer};
\node[ag,left=of wri] (rev) {Reviewer};

\draw[ar] (lit) -- (ida);
\draw[ar] (ida) -- (hyp);
\draw[ar] (hyp) -- (plan);
\draw[ar] (plan) -- (code);
\draw[ar] (code) -- (wri);
\draw[ar] ([yshift=1mm]wri.west) -- ([yshift=1mm]rev.east);
\draw[ar] ([yshift=-1mm]rev.east) -- ([yshift=-1mm]wri.west);
\node[lbl,below=0.5mm of wri] {revise $\rightleftarrows$ review};
\end{tikzpicture}
\caption{The seven-stage pipeline. The Writer and Reviewer form a cycle that terminates
when the Reviewer reports no issues and every quality score clears a threshold, or when an
iteration limit is reached.}
\label{fig:architecture}
\end{figure}

\begin{description}[style=unboxed,leftmargin=0pt,itemsep=0.2em]
  \item[Literature] queries arXiv, Semantic Scholar~\cite{semanticscholar} and
    CORE~\cite{corapi} in parallel, deduplicates the results, and downloads the available PDFs.
  \item[Interdisciplinary Literature] identifies adjacent fields whose methods could inform
    the same problem, searches each of them with the same clients, and returns the merged pool
    together with \emph{bridge insights} tying cross-field work back to the question.
  \item[Hypothesis] synthesises the paper set into exactly three grounded, testable
    hypotheses, and ranks them against one another on feasibility, testability and grounding.
  \item[Experiment Planner] converts each hypothesis into an implementation-ready plan:
    variables, design, data requirements, methods, evaluation criteria and ordered steps,
    with an explicit feasibility judgement.
  \item[Coder] generates a runnable experiment per feasible plan and, where the environment
    permits, executes it and captures structured results.
  \item[Writer] drafts a paper section by section and renders it to
    PDF via reportlab~\cite{reportlab}.
  \item[Reviewer] checks the rendered paper against the same upstream evidence the Writer
    used, and produces structured feedback the Writer can act on.
\end{description}

Every one of the seven is internally a LangGraph \texttt{StateGraph}, not only the top-level
orchestrator. This is a deliberate cost. A graph is heavier than a function call, and for
several of these agents the control flow would fit comfortably in a Python \texttt{for} loop.
The justification is inspection: because each model call and each deterministic check is an
addressable node over checkpointed state, a run can be resumed from the stage at which it
stopped (NFR4) and examined step by step afterwards --- which, given that the user is assumed
unable to read the generated code, is the only remaining means of establishing what happened.

\section{Chaining by data shape}
\label{sec:chaining}

The agents are not coupled through a shared orchestrator object. Each entry point,
\texttt{run\_<name>\_agent()}, accepts a plain dictionary shaped like the \emph{previous
agent's output}, never that agent's internal state, and validates it against the upstream
agent's own \texttt{validate\_output()} before doing anything else (NFR3). A malformed or
outdated file therefore fails immediately with a clear error rather than confusing the model
downstream.

The practical consequence is that every link works identically in process or through disk: a
previous run's JSON can be read from \texttt{outputs/} and passed straight into the next
stage. This is what makes FR10 --- running only selected stages --- a configuration rather
than a separate execution path, and it is what allowed the Interdisciplinary Literature agent
to be inserted into a working pipeline without modifying either of its neighbours. That agent
publishes its merged pool under a \texttt{papers} key, which is precisely the shape the
Hypothesis agent already consumed.

Optionality is preserved additively rather than by changing contracts. The Hypothesis agent's
\texttt{interdisciplinary\_context} and the Planner's \texttt{hypothesis\_ids} both default to
behaving exactly as they did before those parameters existed.

\section{The determinism boundary}
\label{sec:determinism}

NFR1 states that every decision which can be settled mechanically must be settled in code
rather than by model judgement. This is the organising principle of the system, and the
requirements analysis in Chapter~\ref{ch:requirements} is its justification: no participant
would accept unchecked output, and one located the objection precisely in the opacity of model
decisions. The response is not to make the model more persuasive but to reduce the set of
things it is trusted to decide.

Table~\ref{tab:determinism} lists the principal instances. Each was introduced in response to
an observed failure, not anticipated in advance --- a point returned to in
Section~\ref{sec:abandoned} and in Chapter~\ref{ch:bcs}.

\begin{table}[htbp]
\centering\small
\caption{Decisions moved out of the model, and the failure that motivated each.}
\label{tab:determinism}
\begin{tabularx}{\textwidth}{lXX}
\toprule
\textbf{Decision} & \textbf{Left to the model} & \textbf{Now settled in code} \\
\midrule
Is this citation real?
  & Cite the retrieved papers
  & Markers resolve only against the retrieved index; anything else is stripped and reported \\
Was the hypothesis supported?
  & Judge the outcome
  & Computed from the Coder's own execution status and success criteria \\
Are these two papers the same?
  & Judge duplicates
  & One DOI / normalised-title key, shared by both literature agents \\
Which hypothesis wins?
  & Assert a selection
  & Derived from whichever ranking entry holds rank 1; the schema rejects a ranking that is
    not a permutation of $1\ldots3$ \\
Is the data read guarded?
  & Follow the prompt's instruction
  & The generated \texttt{load\_data} is AST-parsed and unguarded reads are rejected \\
Is the paper accurate?
  & Score its own work
  & Citations, metrics and coverage checked by regular expression and dictionary comparison \\
\bottomrule
\end{tabularx}
\end{table}

The last row is the one the literature most directly supports. Agent
Laboratory~\cite{schmidgall2025agentlab} reports that its automated reviewer substantially
overestimates paper quality relative to human reviewers, and Gupta and
Pruthi~\cite{gupta2025glitters} show that model-based novelty checks miss plagiarism entirely.
The Reviewer implemented here therefore re-derives its own copy of the ground truth
independently rather than reading the Writer's summary, and only two things are put to the
model at all: nuance in detecting hallucination, and honesty of framing and tone. Everything
checkable is checked.

The fifth row deserves emphasis because it is the clearest illustration of the principle. The
code-generation prompt had instructed the model for months to guard every data read and fall
back to synthesised data if the file was absent. A production run then generated code that
assumed a \texttt{survey\_data.csv} would be present for a plan whose data requirement was to
collect new data. The instruction was not wrong; nothing checked it. An instruction that is
not verified is a hope.

\section{Graph structure: where concurrency is real, and where it is forbidden}
\label{sec:graphs}

The seven graphs take three shapes, chosen by what the stage actually does rather than by
uniformity.

\paragraph{Fan-out over a fixed set.} The Literature agent queries three independent HTTP APIs
and merges the results; the Reviewer runs a fixed set of independent checks. Both use ordinary
edges, since the branches are known at compile time.

\paragraph{Fan-out over a dynamic set.} The Hypothesis agent batches large paper sets, the
Planner plans each hypothesis independently, and the Interdisciplinary agent searches each
adjacent field. The number of branches is known only at run time, so these use LangGraph's
\texttt{Send} API. One subtlety: the Interdisciplinary agent's conditional edge returns a
plain node name when no adjacent field was identified, because an empty \texttt{Send} list
would otherwise leave the run with no path to its output node.

\paragraph{Cycles, deliberately sequential.} The Coder is the only agent with true cycles, and
the only one where concurrency is actively forbidden. Its per-plan loop and its nested fix
loop are graph cycles rather than Python loops, for the same observability reason as above.
They are not \texttt{Send} fan-outs because the plans are not independent: the cap on SLURM
submissions per run is enforced against a counter that each plan must read after every
preceding plan has updated it. Running them concurrently would let two plans observe the same
pre-submission count and both submit, silently changing real submission-gating behaviour. The
sequencing is a correctness property, not a missed optimisation; the per-plan work is I/O and
model bound, so nothing is lost.

Node-level retry follows the same reasoning. Model- and search-calling nodes carry a retry
policy as an outer net over the retries already beneath them, but nodes with side effects are
excluded explicitly: retrying the Coder's fix loop would re-execute generated code or
re-provision an environment. The policy also uses LangGraph's default exception set, which
does \emph{not} retry \texttt{ValueError} --- a schema failure that survived a repair attempt
is persistent, not flaky.

\section{Transporting generated code}
\label{sec:transport}

Agents that ask the model for short structured fields use a shared helper that invokes the
model, parses JSON, and retries once with a repair prompt. The Coder cannot use it, and the
reason is instructive.

Putting source code inside a JSON string value forces every newline, quote and backslash
through an escaping round-trip that the model must perform by hand. A 3B-active-parameter
model does this unreliably: a literal \texttt{\textbackslash d} in a regular expression, a
real newline where \texttt{\textbackslash n} was required, a trailing backslash that
terminates nothing. Three separate deterministic repairs had accumulated in the codebase to
undo these corruptions before it became clear that the transport itself was the defect.

Code is therefore carried unescaped between delimiters, in a format with no escaping anywhere,
so that a real newline is a real newline and a real backslash is a real backslash:

\begin{lstlisting}[caption={The delimited section format used for all code-bearing responses.},label={lst:sections}]
===BEGIN load_data_function===
def load_data():
    path = DATA_DIR / "corpus.csv"       # a real newline is a real newline
    try:
        return pd.read_csv(path)
    except FileNotFoundError:
        return synthesize_corpus()       # guarded, per check_data_fallback
===END load_data_function===
\end{lstlisting}

The parser mirrors the JSON helper's contract exactly --- one repair retry, then raise --- so a
call site moves between them without other changes. All three deterministic repairs were
deleted when it landed. The section names the prompt asks for are generated from the same list
the parser accepts, so the requested shape cannot drift from the accepted shape.

\section{Making generated experiments work}
\label{sec:coder}

The Coder agent is the largest component in the system and the one where aim A5 --- that
experiments recover from their own failures, because the intended user cannot debug them ---
is discharged.

\subsection{The template splice}

An experiment is a single \texttt{run.py} rendered from a fixed template. The metadata block
and the orchestration footer --- timing, exception handling, and writing
\texttt{results.json} --- are not model-generated. Only \texttt{load\_data},
\texttt{build\_model}, \texttt{run\_experiment} and \texttt{evaluate}, plus imports and
configuration, are, and they are spliced into the template by plain string replacement rather
than \texttt{str.format()}, because generated code routinely contains literal braces that
would break format-string substitution.

This replaced an earlier design in which the model wrote four separate files. That design
required the model to observe a cross-file calling convention consistently, and it got it
slightly wrong on roughly one experiment in every batch. Removing the wiring from the model's
responsibilities removed the entire failure mode rather than mitigating it.

\subsection{Classifying a failure before responding to it}

Every runtime failure was originally recorded as \texttt{run\_experiment} --- the stage that
broke, never what broke --- and every such failure took the same route: regenerate the code and
try again. That is correct for a logic bug and useless for anything else. A production run
demonstrates the cost: three fix attempts, three regenerations, and
\texttt{ModuleNotFoundError: No module named 'pandas'} returned verbatim each time. Rewriting
the source could not install anything.

Failures are now classified by kind before anything decides what to do about them, as shown in
Table~\ref{tab:failures}. Three kinds never reach the model at all.

\begin{table}[htbp]
\centering\small
\caption{Failure classification and routing. The first three routes spend no fix attempt,
because none represents a defect in the generated source.}
\label{tab:failures}
\begin{tabularx}{\textwidth}{llX}
\toprule
\textbf{Kind} & \textbf{Route} & \textbf{Action} \\
\midrule
Missing dependency & environment & Install the distribution, verify the import now resolves,
  re-run the \emph{unchanged} code \\
Resource limit & downscale & Halve the code's own cost parameters and re-run \\
Removed API & regenerate & Rewrite the deleted call mechanically \\
Dead import & regenerate & Regenerate carrying the successor API, because installing
  cannot help \\
Missing system library & terminal & Report; pip cannot supply it \\
Logic defect & regenerate & The fix loop proper \\
\bottomrule
\end{tabularx}
\end{table}

Two distinctions in that table are subtle and were learned the hard way. An import that pip
\emph{can} satisfy (\texttt{sklearn} $\rightarrow$ \texttt{scikit-learn}) is an alias, and
installing works. An import that no package provides (\texttt{pymc3}, \texttt{pystan}) is a
succession, and installing the successor reports success while changing nothing --- the
identical error returns on the re-run and the loop makes no progress. These route in opposite
directions. Similarly, the removed-API case is regenerated \emph{carrying the replacement},
because the traceback states what broke but not what to write instead; this was added
specifically because the fix prompt already named the correct replacement and the model
reproduced the identical deleted call regardless. Asking again is the thing that does not work.

\subsection{Making a failed attempt cheap before making it rare}

The first execution of any generated experiment is a \emph{smoke run}: the same code with every
known cost parameter pinned to its floor, under a short timeout rather than the complexity
timeout. A defect anywhere in the program therefore surfaces in seconds rather than after a
full-length run --- and so does each subsequent round of the fix loop.

The asymmetry is deliberate and load-bearing. A smoke run can end an attempt early; it can
never let one through. Anything that passes is still executed at full size, and the smoke
run's own results file is deleted rather than read. A failure that the shrinking could itself
have caused is re-run at full size rather than reported.

A related check operates one step earlier. A program that parses is not yet a program that can
run, so the generated source is checked with pyflakes for undefined names before an environment
is provisioned. The same defect previously cost a full environment provision, an execution, and
a fix attempt.

\subsection{Two budgets, because two things can go wrong}

The fix loop is bounded, but not by a single counter. A response that is \emph{malformed or
hollow} rather than wrong --- unparseable, missing a section, a function under the wrong name,
a body consisting of \texttt{pass} --- rendered nothing, provisioned nothing and executed
nothing. It draws on a separate structural-retry budget, leaving the fix budget for the defect
underneath. The failure that forced this split was a run which spent all three fix attempts on
the model omitting a section, never reaching the code at all.

A compile error deliberately remains on the fix budget: a syntax error is a real defect in a
real answer. Exhausting the structural budget ends the plan rather than falling through to the
fix budget, because the identical-failure stop condition cannot bound that fall-through --- an
unparseable response's summary quotes the model's own varying output, so repeated failures do
not compare equal.

\subsection{Confining a regeneration}

The fix prompt instructs the model to keep whatever already worked. That is an instruction, not
a guarantee. A production trace records attempt~1 correctly fixing a \texttt{ModuleNotFoundError}
while attempts~2 and~3 each introduced a fresh syntax error in a region they were never asked
to touch.

A check that can \emph{localise} its failure --- to a section it parsed independently, or to a
line the template's span map resolves back to one --- therefore requests only those sections and
reuses the rest verbatim. And because a fix loop is not monotonic, giving up reports the attempt
that got \emph{furthest}, ranked by the same ordering that defines whether a regeneration made
progress, rather than whichever attempt happened to be last. Ties go to the final attempt, so a
run that never regressed is unaffected.

\section{Provenance, and the withheld verdict}
\label{sec:provenance}

Aim A2 requires that results computed on synthesised inputs withhold a conclusion. Every input
a plan declares is resolved to a real source or a surrogate, and any surrogate sets the
experiment's success criterion to the string \texttt{"unknown"}, which the Writer maps to
\emph{inconclusive}.

The alternative --- returning \texttt{False} --- would map to \emph{refuted}, meaning the paper
would claim to have refuted a hypothesis on the basis of numbers computed over invented inputs.
That is the specific outcome the mechanism exists to prevent. The metrics are still reported
and the model's own claim is retained under a separate field; only the verdict is withheld.

Two deterministic passes run afterwards, in order. The first downgrades a declared download the
code never actually fetched --- the over-claiming direction, in which a plan that quietly
synthesised data would otherwise still earn an interpretable verdict. The second drops a
\emph{phantom} surrogate: a requirement no source could be found for, which a real input the
code demonstrably reads already answers under a different name. Only unresolved entries are
superseded; a restricted or credentialed source names real data that specifically was not
obtained, and nothing else answers it.

\section{Grounding the write-up}
\label{sec:writing}

The Writer drafts section by section, each call seeing only the slice of upstream data that
section needs, so the full pipeline output is never forced into a single context window.

Citations use markers of the form \texttt{[[cite:ID]]} and \texttt{[[citet:ID]]}, restricted to
identifiers actually present in the retrieved paper index --- the same identifier scheme the
Hypothesis agent used, so citations line up with the paper ids already referenced by the
hypotheses. Resolution is necessarily two-phase, because author-year citation cannot format any
single marker until the whole cited set is known: two papers may share a first-author surname
and year and require a \emph{2020a}/\emph{2020b} suffix. Sections are therefore scanned first,
disambiguation and reference ordering are fixed once, and only then is each marker replaced.

A marker naming an identifier that is not in the index is never trusted into the output. It is
stripped and recorded for review, so a fabricated citation cannot reach the rendered PDF. The
Writer and Reviewer are given the \emph{same} paper set --- the merged in-domain and cross-field
pool --- because otherwise the Reviewer would flag citations the Writer was entitled to make.

Rendering is via reportlab rather than LaTeX. This deviates from the project proposal, and the
justification is operational: reportlab is pure Python, so the Writer runs unmodified on the
HPC facility with no TeX installation, and the proposal's own risk register had rated
``generated LaTeX reports fail to compile'' as high-likelihood. The risk was removed rather
than mitigated.

\section{The web application}
\label{sec:webapp}

The interface (FR11) reads the pipeline and never modifies it. Progress comes from streaming
node-level updates out of the graph, and the display layer reduces each node's own delta to
display fields; no agent logic is reimplemented for it.

It runs as three processes, not one. The server never executes pipeline work: starting a run
spawns a separate runner process, and the server thereafter only reads files that process
writes. This is what makes a run cancellable at all --- a thread inside a blocking graph
invocation cannot be interrupted, whereas a subprocess accepts a signal --- and it additionally
isolates the server from a crash inside generated code, permits the server to be restarted
mid-run, and gives each run its own experiment directory, which is the one output location the
graph itself cannot route per run. Each run is a self-contained directory, so the server's
entire view of a run is the files on disk.

\section{Deployment}
\label{sec:deployment}

The pipeline is designed to run on the university's HPC facility (A4). Model weights are served
by vLLM~\cite{kwon2023vllm} inside a container, and the pipeline runs on the same allocated
node so that model traffic never leaves it. Because the agents are CPU-bound and the GPU nodes
have many cores, co-locating costs nothing and avoids both an SSH tunnel and any long-running
process on a login node.

Aim A6 is addressed by durable checkpointing. All eight compiled graphs obtain their
checkpointer from one factory, so the whole pipeline switches between in-process and durable
storage through a single environment variable rather than eight edits. Durability is what makes
resumption possible, and both unattended entry points use it: a failed run is offered a resume
control in the interface, and a batch sweep records each question's thread identifier in a
manifest so that a pre-empted job resumes mid-question rather than restarting it. The resume
control is hidden entirely under in-process checkpointing, rather than offering a restart that
describes itself as a resume.

A subtlety in the batch path is worth recording. Its original per-question identifiers were
freshly generated each run and so were, by construction, identifiers no later process could ask
for. Persisting them was necessary before resumption could work at all --- and a recorded
identifier is reused only when the question at that index still matches, because indices are
positions in a file that people edit between sweeps, and resuming the wrong thread would
silently graft half a run about one question onto another.

\section{Approaches tried and abandoned}
\label{sec:abandoned}

Four are worth recording, beyond those already described.

\paragraph{An in-process model backend.} An alternative backend loading weights directly rather
than through an HTTP server was built and later removed, once the pipeline settled on a single
served model. It had allowed smaller models on hardware too small for the target model, but it
doubled the configuration surface for a path nothing used.

\paragraph{A change of model.} The pipeline originally targeted a reasoning model that emitted
an explicit thinking trace. Moving to a coding-specialised model removed the need for that
trace, but the stripping logic was retained rather than deleted, so that pointing the
configuration back at a reasoning model requires no code change.

\paragraph{Resolving symbolic links.} An apparently innocuous call to resolve a virtual
environment's interpreter path to its canonical form followed the symbolic link to the base
interpreter and discarded the environment entirely. Two cluster jobs were lost to it before the
cause was found. Absolute paths were required; canonical paths were not.

\paragraph{Trusting an instruction.} Discussed in Section~\ref{sec:determinism}, and worth
stating as a general lesson rather than a single bug: several mechanisms in this chapter exist
because a prompt had asked for behaviour that nothing verified, and the behaviour did not
occur. Every such instruction that could be checked in code eventually was.
